\section{Fetch API dan Permintaan HTTP dari Browser}

\textbf{Fetch API} adalah interface modern untuk melakukan permintaan HTTP dari JavaScript di browser. \code{fetch(url)} mengembalikan Promise yang resolve ke objek \code{Response}. Metode \code{.json()} pada Response juga mengembalikan Promise. Fetch mendukung semua metode HTTP (GET, POST, PUT, DELETE) melalui objek opsi \cite{jsInfo}.

\begin{contoh}
\textbf{Studi Kasus: Menampilkan Data Produk dari API}

Aplikasi yang mengambil daftar produk dari REST API dan menampilkannya di halaman:
\begin{enumerate}
  \item Tampilkan loading indicator
  \item \code{const res = await fetch("/api/products")}
  \item \code{const products = await res.json()}
  \item Loop products; buat elemen DOM untuk setiap produk
  \item Sembunyikan loading; tampilkan daftar produk
  \item Tangani error: tampilkan pesan jika network gagal
\end{enumerate}
\end{contoh}

\begin{webcode}[language=JavaScript, caption={Fetch API: GET dan POST}]
// GET request
async function getProducts() {
  const res = await fetch("/api/products");
  if (!res.ok) throw new Error(`HTTP ${res.status}`);
  return res.json();
}

// POST request (kirim data)
async function createProduct(data) {
  const res = await fetch("/api/products", {
    method: "POST",
    headers: { "Content-Type": "application/json" },
    body: JSON.stringify(data)
  });
  if (!res.ok) throw new Error(`HTTP ${res.status}`);
  return res.json();
}
\end{webcode}

\begin{catatan}
\textbf{fetch() Tidak Reject pada HTTP Error.} Berbeda dari XMLHttpRequest, \code{fetch()} hanya reject pada network error (tidak ada koneksi). Response 404 atau 500 tetap resolve---periksa \code{res.ok} atau \code{res.status} untuk mendeteksi error HTTP. Ini adalah ``gotcha'' paling umum saat menggunakan Fetch API \cite{mdnJsGuide}.
\end{catatan}

